% 由 scripts/make_paper_tables.py 生成（同预算同题集配对比较），请勿手改
\begin{table}[t]
  \centering
  \caption{Prompt scaffolding as a substitute for capability. The same
           150 problems are evaluated with a matched budget (2048 new tokens) under
           two prompts: standard chain-of-thought, and a structured prompt that additionally
           asks the model to write down the knowns and unknowns explicitly. Accuracy does not
           improve significantly, so the instability documented in
           Table~\ref{tab:sensitivity} is not a missing-scaffold problem.}
  \label{tab:scaffold}
  \small
  \begin{tabular}{@{}l c c c c@{}}
    \toprule
    \textbf{Prompt} & \textbf{Accuracy} & \textbf{95\% CI} & \textbf{No \texttt{ANSWER:}} &
      \textbf{Output chars (mean / max)} \\
    \midrule
    standard cot & 24.0\% & [17.3, 31.3] & 71\% & 2377 / 4224 \\
    structured (knowns/unknowns) & 26.0\% & [19.3, 33.3] & 90\% & 2373 / 4639 \\
    \midrule
    \multicolumn{4}{@{}l}{paired difference: +2.0\,pp, McNemar $p=0.700$} & \\
    \midrule
    \multicolumn{5}{@{}l}{\emph{Accuracy by output-format compliance (both prompts pooled)}} \\
    \quad outputs with \texttt{ANSWER:} line (n=58) & 20.7\% & -- & -- & -- \\
    \quad outputs without it (n=242) & 26.0\% & -- & -- & -- \\
    \bottomrule
  \end{tabular}
  \vspace{2pt}
  \parbox{\linewidth}{\footnotesize Upper block: scaffolding changes accuracy by less than
  one percentage point and is far from significant. Lower block: most outputs ignore the
  requested \texttt{ANSWER:} format and are graded through the fallback extractor;
  accuracy is similar with and without the marker, so the fallback does not visibly bias
  the headline numbers, while output lengths stay well below the token budget --- a
  format-compliance failure rather than truncation.}
\end{table}
